% ---------------------------------------------------------------------
%  Chapter 6 : The Field-Theory Core --- Action, Taylor Expansion, Vertices
% ---------------------------------------------------------------------
\chapter{The Field-Theory Core: Action, Taylor Expansion, Vertices}\label{ch:fieldtheory-core}

This is the first of the machinery chapters, and it opens the very first of the
seven phases you watched scroll past in the quickstart trace: \code{[1/7]
FieldTheory.expand}. Everything downstream --- the propagator, the mean-field
solve, the diagram enumeration, the loop integrals --- consumes the output of
this one stage. If you understand nothing else about the internals, understand
this: the field-theory core is the machine that turns a \emph{declarative model}
(a Python dictionary of strings, lists, and small functions) into a
\emph{classified symbolic action} --- a dictionary of polynomial pieces sorted by
how many response legs and how many physical legs each piece carries. The free
propagator, the noise sources, and the interaction vertices all fall out of that
single classification.

The subsystem lives in four files, all under \file{msrjd/core/}:

\begin{itemize}
  \item \file{field\_theory.py} --- the \code{FieldTheory} expander itself: it
        builds the symbolic namespace, evaluates the action, Taylor-expands every
        nonlinearity, classifies the result by ``bigrade,'' and verifies the
        mean-field sector vanishes.
  \item \file{vertices.py} --- decomposes each classified sector into individual
        typed records (\code{VertexType}, \code{SourceType}, and their
        kernel-carrying subclasses) that the diagram machinery can build with.
  \item \file{convolution.py} --- the formal \code{Conv(kernel, field)} operator
        and the rules that resolve it, so a synaptic kernel times a field means
        \emph{convolution in time}, not ordinary multiplication.
  \item \file{serialize.py} --- saves an expanded theory to disk (so the
        expensive Taylor step is not repeated) and re-imports a model on demand.
\end{itemize}

We will work two examples in parallel. The running ``toy'' is the quartic
Ornstein--Uhlenbeck model from Chapter~\ref{ch:quickstart} --- one physical field,
white noise, a single cubic vertex --- because it is small enough that you can
hold the entire expanded action in your head. The running ``real'' example is the
quadratic Hawkes model shipped in \file{models/hawkes\_quad\_expg.py} --- two
populations, a synaptic convolution kernel, a formal transfer function --- because
it exercises every feature the subsystem has. Reading that model file alongside
this chapter is recommended; it is the canonical shape of the ``model dict'' the
core consumes.

\begin{note}
A vocabulary reminder from the physics, since the code leans on it constantly. In
the \msrjd{} construction every \term{physical field} (the thing that actually
fluctuates --- a voltage $v$, a spike rate $n$, our $x$) is paired with a
conjugate \term{response field}, written with a tilde ($\tilde v$, $\tilde n$,
$\tilde x$) and called the ``tilde'' or ``hat'' field. In diagram language a
physical leg is an ``in-leg'' and a response leg is an ``out-leg.'' The entire
classification scheme of this chapter is built on counting these two kinds of
legs separately.
\end{note}

\section{The MSR-JD action as data}

The single most important design choice in this subsystem --- the one that makes
it model-independent --- is that \emph{a model is data, not code}. You do not
write a new expander for each physics problem. You hand the one expander a Python
dictionary describing your model, and it does the rest. So before any algebra, we
need to understand the shape of that dictionary and the object the framework
builds from it.

\subsection{The model dictionary}

A model is a \code{dict}. Its values are a mix of plain data (strings, numbers,
lists of small spec dictionaries) and \emph{lambdas} --- small functions that,
given a namespace object, return a symbolic expression. The keys the core reads
are summarized in Table~\ref{tab:modeldict}.

\begin{table}[h]
\centering
\caption{The keys of a model dict that the field-theory core reads. ``lambda
\code{ns}$\to$\,\dots'' means a function taking the namespace object and returning
the indicated thing. Optional keys are noted.}
\label{tab:modeldict}
\small
\begin{tabular}{@{}lll@{}}
\toprule
key & type & meaning \\
\midrule
\code{name}              & str  & human-readable model name \\
\code{index\_sets}       & dict & named index ranges; first is the primary index \\
\code{populations}       & list & per-population \code{\{name, size\}} (optional) \\
\code{response\_fields}  & list & specs for the tilde fields \\
\code{physical\_fields}  & list & specs for the fluctuation fields \\
\code{parameters}        & list & \code{\{name, indexed\_by, domain, default, \dots\}} \\
\code{kernels}           & list & \code{\{name, sage\_name, indexed\_by, \dots\}} \\
\code{operators}         & list & \code{\{name, sage\_name\}} (e.g.\ \code{Dt}) \\
\code{functions}         & list & nonlinear fns: \code{expression} or \code{taylor\_coeffs} \\
\code{mf\_substitutions} & list & \code{\{name, value: lambda ns$\to$SR\}} \\
\code{mf\_bg\_conditions}& lambda & saddle substitutions for sector verification \\
\code{specializations}   & lambda & rename abstract derivative symbols to values \\
\code{correlated\_noises}& dict & colored-noise cumulant declarations (optional) \\
\code{kernel\_ft\_image} & lambda & $\{$kernel\,$\mapsto\hat g(\omega)\}$ (Fourier image) \\
\code{kernel\_td\_image} & lambda & $\{$kernel\,$\mapsto g(\tau)\}$ (time-domain image) \\
\code{action}            & lambda & \textbf{the action} \\
\bottomrule
\end{tabular}
\end{table}

You will rarely author this dict by hand. The theory builders of
Chapter~\ref{ch:theory-spec} generate it for you from a much friendlier fluent
API (recall
\code{TemporalTheoryBuilder('OU Quartic')\allowbreak.physical\_field('x')\dots}).
But the core does not know about the builders; it only sees the dict. So
everything in this chapter is phrased in terms of these keys.

The one key that does the real work is \code{action}: a function that, handed the
namespace, returns one big symbolic expression for $S$. For the quartic OU model
the builder produces an action equivalent to
\[
  S \;=\; \sum_{i}\Bigl[\tilde x_i\,\bigl((\partial_t + \mu)\,x_i + \varepsilon\,x_i^{3}\bigr)
          \;-\; D\,\tilde x_i^{2}\Bigr],
\]
which is exactly the action text \code{xt[i]*((Dt+mu)*x[i] + eps*x[i]\^{}3) - D*xt[i]\^{}2}
you saw in the quickstart theory file. For the Hawkes model the action authored at
\file{hawkes\_quad\_expg.py:236} is the richer
\[
  S \;=\; \sum_i\Bigl[\tilde n_i\,(\dot{n}^{*}_i + \delta n_i)
        \;-\;(\ee^{\tilde n_i}-1)\,\phi_i(\delta v_i)
        \;+\;\tilde v_i\bigl[(\tau\partial_t+1)\delta v_i + v^{*}_i - E_i
              - \textstyle\sum_j w_{ij}\,g\!\cdot\!(n^{*}_j+\delta n_j)\bigr]\Bigr],
\]
with response fields $\tilde n_i,\tilde v_i$, fluctuation fields
$\delta n_i,\delta v_i$, a synaptic kernel $g$ (a convolution), and a nonlinear
transfer function $\phi_i(v)=a\,v^2$.

\subsection{The namespace object}

The action lambda takes one argument, conventionally called \code{ns}, and writes
the action in terms of attributes of \code{ns}: \code{ns.xt[i]}, \code{ns.x[i]},
\code{ns.mu}, \code{ns.Dt}, and so on. That \code{ns} object is an instance of a
deliberately empty class:

\begin{lstlisting}
# msrjd/core/field_theory.py:185
class _Namespace:
    pass
\end{lstlisting}

It is empty because it is a blank slate that \code{FieldTheory.\_build\_namespace}
fills, one attribute at a time, with a symbolic variable for every field, every
parameter, every kernel, every operator, and every function the model declares
(we walk that construction in detail in \S\ref{sec:buildns}). The action lambda
never has to import anything model-specific; it just reaches into the namespace it
was handed. This indirection is the whole trick that lets one expander serve every
model.

\begin{gotcha}
\code{\_Namespace} also carries a pile of private bookkeeping attributes whose
names start with an underscore: \code{ns.\_all\_field\_sr\_vars} (the list of every
field symbol), \code{ns.\_ring\_var\_names} (the ring generator names, tilde
first), \code{ns.\_deriv\_rename\_subs} (the formal-derivative rename map), and
more. These are not for the model author to touch; they are how the expander
passes state from \code{\_build\_namespace} to the rest of \code{expand}. The
namespace is also \emph{unpicklable} --- it holds inner closures --- which becomes
important in \S\ref{sec:pickle}.
\end{gotcha}

\section{SageMath in one page}

Every field, parameter, kernel, and operator in this subsystem is a SageMath
symbolic object. If you have never used Sage, this section is the minimum you need
to read the code; everything here recurs throughout the manual.

\term{SageMath} is a large open-source mathematics system built on Python. It
bundles many computer-algebra engines (Maxima, Pari, Singular, FLINT, GMP, SymPy,
\dots) behind one Python API. This subsystem uses exactly two pieces of it.

\paragraph{The Symbolic Ring \code{SR}.} \code{SR} is Sage's universal ring of
symbolic expressions --- think ``anything you could write on a blackboard'':
\code{a*x\^{}2 + sin(t) + exp(nt)}. \code{SR.var('a')} creates a symbolic variable
named \code{a}. Wrapping a value in \code{SR(...)} coerces it into the symbolic
ring, so \code{SR(0)} is the symbolic zero. Crucially, an \code{SR} variable can
carry an \emph{assumption} (e.g.\ \code{domain='positive'}, which tells Sage the
symbol is a positive real) and a \emph{LaTeX display name} (\code{latex\_name=...},
so a notebook renders $\mu$ instead of \code{mu}).

\paragraph{\code{PolynomialRing(SR, names)}.} This builds a polynomial ring
\emph{whose coefficients live in \code{SR}} and whose \emph{generators} are the
named field variables. So \code{PolynomialRing(SR, ['xt','dx'])} represents the
action as a genuine polynomial in the two field generators \code{xt} and
\code{dx}, while the parameters ($\mu$, $\varepsilon$, $D$) live symbolically
\emph{in the coefficients}. This separation is the linchpin of the whole chapter:
a polynomial-ring element exposes its monomials through \code{.dict()}, a map
\code{\{exponent\_vector: coefficient\}}, and the exponent vector immediately tells
you how many of each field a monomial contains.

The exact imports the expander uses make the surface concrete:

\begin{lstlisting}
# field_theory.py:26-29
from sage.all import (
    SR, PolynomialRing, factorial, QQ, latex, LatexExpr,
    diff, function, exp, dirac_delta, heaviside, integrate, oo, I, pi, taylor
)
\end{lstlisting}

A quick gloss of the less obvious names: \code{factorial} and \code{QQ} (Sage's
field of \emph{exact} rationals) build Taylor coefficients $1/k!$ without ever
touching floating point; \code{diff} is symbolic differentiation; \code{function}
builds a \emph{formal} (unevaluated) symbolic function (\S\ref{sec:formal});
\code{dirac\_delta} and \code{heaviside} are $\delta(t)$ and the step function;
\code{integrate} is symbolic integration; \code{oo} is $\infty$, \code{I} is
$\sqrt{-1}$, \code{pi} is $\pi$; and \code{taylor} is the multivariate Taylor
expander (\S\ref{sec:taylor}).

Two Sage behaviors trip up newcomers so hard that the code comments call them out
by name. They are worth a callout box of their own, because getting either wrong
\emph{silently corrupts the answer}.

\begin{gotcha}
\textbf{A Sage sum is \code{add\_vararg}, not \code{operator.add}.} When you ask a
Sage expression for its top-level operator with \code{expr.operator()}, a sum does
\emph{not} report Python's \code{operator.add}; it reports Sage's own internal
n-ary callable \code{add\_vararg}. So the code never compares operators by
identity. Everywhere it needs to ask ``is this a sum?'' it does a name-based test:
\begin{lstlisting}
# field_theory.py:266-267  (inside _sr_to_ring)
_op_obj = expanded.operator()
if _op_obj is not None and 'add' in getattr(_op_obj, '__name__', '').lower():
    summands = expanded.operands()
\end{lstlisting}
The same pattern appears in \file{convolution.py:254}. Getting this wrong does not
raise an error --- it just fails to recognize the sum, and terms are silently
dropped.
\end{gotcha}

\begin{gotcha}
\textbf{\code{expr != 0} is unreliable for symbolic zero.} Under Sage, when an
expression contains an unbound real parameter with no positivity assumption (e.g.\
a symbol declared \code{domain='real'}), the comparison \code{expr != 0} can return
\code{False} --- Sage cannot decide the sign and refuses to commit. The correct
\emph{syntactic} zero test is \code{.is\_trivial\_zero()}, which is \code{True}
exactly when the expression is literally zero, regardless of any parameter
assumptions. The cumulant injector relies on it (\file{field\_theory.py:490}),
with the reason spelled out in the comment right above the call.
\end{gotcha}

\section{From action to polynomial: the \texttt{expand()} pipeline}

With the vocabulary in place we can read the heart of the subsystem:
\code{FieldTheory.expand} at \file{field\_theory.py:814}. It is a single method
that performs an eight-step transform. We list the steps first, then spend the
following sections opening up the ones that carry real subtlety (Taylor expansion,
convolution, bigrade classification, mean-field verification). Here is the
skeleton, lightly trimmed:

\begin{lstlisting}
# field_theory.py:814  (FieldTheory.expand, abridged)
def expand(self) -> None:
    ns, R, n_tilde = self._build_namespace()           # 1. build namespace
    self._ns, self._R, self._n_tilde = ns, R, n_tilde

    S_sr = SR(self.model['action'](ns))                # 2. evaluate the action

    # 3. resolve Conv(kernel, field) atoms
    S_sr = reduce_conv_in_action(S_sr, fluct_vars,
               taylor_order=self.taylor_order, attachments_out=attachments)
    ns._kernel_attachments = attachments

    # 4. inject correlated-noise cumulant series
    S_sr = S_sr + _build_cumulant_action(ns, self.model)

    # 5. multivariate Taylor to total degree taylor_order
    if ns._all_field_sr_vars:
        S_sr = taylor(S_sr, *[(v, 0) for v in ns._all_field_sr_vars],
                      self.taylor_order)

    # 6. rename formal-derivative symbols  D[0](phi_1)(0) -> phi1_1
    S_sr = S_sr.subs(ns._deriv_rename_subs)

    # 7. apply model specializations  phi1_i -> 2 a vstar_i, ...
    S_sr = S_sr.subs(self.model['specializations'](ns))

    # 8. coerce to the polynomial ring and classify BEFORE MF subs
    S_poly = _sr_to_ring(S_sr.expand(), R, ns._ring_var_names)
    by_tp  = _collect_bigrade(S_poly, n_tilde)
    # ... then verify+zero the mean-field sector, rebuild S_raw ...
    self._by_tp = by_tp
\end{lstlisting}

Read top to bottom:

\begin{enumerate}
  \item \textbf{Build the namespace.} \code{\_build\_namespace} returns the triple
        \code{(ns, R, n\_tilde)}: the populated namespace, the polynomial ring, and
        the number of response-field generators. Detailed in \S\ref{sec:buildns}.
  \item \textbf{Evaluate the action.} \code{self.model['action'](ns)} calls the
        model's action lambda with the namespace; the result is one large \code{SR}
        expression. Wrapping it in \code{SR(...)} guarantees it is in the symbolic
        ring even if the lambda returned a bare integer.
  \item \textbf{Resolve convolutions.} Every \code{Conv(g, field)} atom is rewritten
        by \code{reduce\_conv\_in\_action} (\S\ref{sec:conv}). The whole step is
        wrapped in a \code{try/except (AttributeError, TypeError)} so a model with
        no convolutions passes through untouched.
  \item \textbf{Inject correlated-noise cumulants.} If the model declares colored
        noise, \code{\_build\_cumulant\_action} appends a cumulant series to the
        action; otherwise it returns \code{SR(0)}, a true no-op (\S\ref{sec:cumulant}).
  \item \textbf{Taylor-expand.} Every nonlinear function of the fields becomes a
        finite polynomial, truncated at total degree \code{taylor\_order}
        (\S\ref{sec:taylor}).
  \item \textbf{Rename derivative symbols.} Sage produces ugly atoms like
        \code{D[0](phi\_1)(0)} when it differentiates a formal function; these are
        renamed to clean polynomial-friendly names like \code{phi1\_1}
        (\S\ref{sec:formal}).
  \item \textbf{Apply specializations.} The model substitutes concrete values for
        those abstract derivative symbols (for $\phi=a v^2$, $\phi''\to 2a$). These
        are pure renames, safe to apply to every sector.
  \item \textbf{Coerce and classify.} The \code{SR} expression is converted into a
        polynomial-ring element by \code{\_sr\_to\_ring}, then split into bigrade
        sectors by \code{\_collect\_bigrade} (\S\ref{sec:bigrade}). Only \emph{after}
        this does the mean-field sector get verified and zeroed (\S\ref{sec:mf}).
\end{enumerate}

\begin{gotcha}
The docstring of \code{expand} (\file{field\_theory.py:818}) still describes step~5
as ``sequential \code{taylor()}.'' The code at \file{field\_theory.py:894} actually
does a single one-shot multivariate \code{taylor}. Trust the code: the expansion is
multivariate, for the reason given in \S\ref{sec:taylor}.
\end{gotcha}

The order of these steps is not arbitrary. The most important ordering constraint
--- classify \emph{before} applying mean-field substitutions --- is the subject of
\S\ref{sec:mf}, and the comment at \file{field\_theory.py:913--922} is worth
reading in the source.

\subsection{Building the namespace}\label{sec:buildns}

\code{FieldTheory.\_build\_namespace} (\file{field\_theory.py:1041}) is the longest
method in the file. You do not need every line, but you should know the order in
which it populates \code{ns}, because that order fixes the ring generator ordering
that the entire classification depends on.

\begin{enumerate}
  \item \textbf{Index sets and populations.} Each entry of \code{model['index\_sets']}
        is bound onto \code{ns}; the first is the ``primary'' index. For each
        \code{model['populations']} entry it builds a $0\!\dots\!(\text{size}-1)$
        index list, so action text can write \code{for i in pop}. Legacy theories
        with no populations fall back to a single flat \code{pop}.
  \item \textbf{Field variables.} For each \code{response\_fields} and
        \code{physical\_fields} spec it creates the \code{SR} symbols: an indexed
        field becomes a list \code{[SR.var(f"\{fname\}\{i+1\}", latex\_name=...)]}
        (so field \code{x} in a population of one becomes the symbol \code{x1}); a
        scalar field becomes a single symbol. It accumulates the tilde symbols and
        the physical symbols separately and stashes them as
        \code{ns.\_tilde\_sr\_vars}, \code{ns.\_phys\_sr\_vars}, and the union
        \code{ns.\_all\_field\_sr\_vars}.
  \item \textbf{The polynomial ring.} It forms
        \code{ring\_var\_names = tilde\_names + phys\_names} --- \emph{tilde names
        first} --- sets \code{n\_tilde = len(tilde\_names)}, and builds
        \code{R = PolynomialRing(SR, ring\_var\_names)}. This is why every later
        exponent-vector split at index \code{n\_tilde} cleanly separates response
        from physical legs.
  \item \textbf{Parameters.} Scalar parameters get one symbol (with a \code{domain}
        assumption if declared); vector/matrix parameters get a flat list of symbols
        sized by population.
  \item \textbf{Kernels and operators.} A scalar kernel becomes one symbol named by
        its \code{sage\_name}; a vector kernel becomes \code{z\_<name>\_<i+1>}; a
        matrix kernel becomes a list-of-lists \code{z\_<name>\_<i+1>\_<j+1>}.
        Operators (like the time derivative \code{Dt}) become one symbol each. It
        also defines \code{ns.delta\_D = SR.var('z\_delta', \dots)} and a primed
        variant. The leading \code{z} in these internal names is deliberate: \code{z}
        sorts \emph{after} \code{phi}/\code{tau} in any Sage product, giving a
        canonical, reproducible factor ordering.
  \item \textbf{Nonlinear functions.} Each \code{model['functions']} spec installs
        a callable on \code{ns} and registers its derivative-rename entries. There
        are two authoring paths --- an \code{'expression'} path and a legacy
        \code{'taylor\_coeffs'} path --- detailed in \S\ref{sec:formal}.
  \item \textbf{Mean-field substitutions.} For each \code{model['mf\_substitutions']}
        it does \code{setattr(ns, sub['name'], sub['value'](ns))}, computing the
        value once at build time (this is how a saddle quantity like \code{vstar}
        gets onto the namespace so the action lambda can reference it).
\end{enumerate}

The method returns \code{(ns, R, n\_tilde)}. The take-away to carry forward: the
ring generators are ordered \emph{all tilde fields, then all physical fields}, and
\code{n\_tilde} is the boundary between the two halves.

\section{Taylor expansion and formal-function renaming}\label{sec:taylor}

A field theory's interactions live in the nonlinear terms of its action: the
$\varepsilon x^3$ of the OU model, the $\ee^{\tilde n}-1$ and the $\phi(v)$ of the
Hawkes model. To compute Feynman diagrams these must be replaced by finite
polynomials. That is what step~5 of \code{expand} does.

\subsection{Why total-degree truncation}

For a single-argument function expanded around the saddle ($\delta\phi=0$), the
Taylor series truncated at order $N$ is the familiar
\[
  f(\delta\phi) \;=\; \sum_{k=0}^{N} \frac{f^{(k)}(0)}{k!}\,\delta\phi^{\,k}.
\]
For a multi-argument function the framework uses the multivariate Taylor series
truncated at \emph{total} degree $N$:
\[
  f(x) \;=\; \sum_{|\alpha|\le N} \frac{1}{\alpha!}\,(\partial^{\alpha} f)(0)\,x^{\alpha},
\]
where $\alpha=(k_1,\dots,k_n)$ is a multi-index, $|\alpha|=\sum_j k_j$,
$\alpha!=\prod_j k_j!$, and $x^{\alpha}=\prod_j x_j^{k_j}$.

Truncating at \emph{total} degree (not per-variable degree) is the physically
correct choice, and it is worth being explicit about why. The total degree of a
monomial in the fields equals the number of legs of the diagrammatic vertex it
produces. So ``Taylor order $N$'' means exactly ``keep every vertex with up to $N$
legs.'' This is the single knob \code{taylor\_order} (default $4$), and it is the
same knob the quickstart's \code{[1/7]} trace printed as \code{taylor\_order=...}.

The code realizes this with one Sage call:

\begin{lstlisting}
# field_theory.py:894
S_sr = taylor(
    S_sr,
    *[(v, 0) for v in ns._all_field_sr_vars],
    self.taylor_order,
)
\end{lstlisting}

Sage's \code{taylor(expr, (v1,0), (v2,0), ..., N)} expands \code{expr} jointly in
all listed variables, each about $0$, to total order $N$. The \code{*[(v, 0) for
...]} splat hands it one \code{(variable, 0)} pair per field.

\begin{gotcha}
This is a \emph{one-shot multivariate} Taylor, and it must be. An earlier version
of the code chained single-variable \code{.taylor()} calls in a loop. That is wrong
for multi-argument formal functions: chained \code{.taylor()} does not compose at
non-zero expansion points, because the outer call freezes the inner-argument
substructure of the formal function. The comment at \file{field\_theory.py:884--892}
documents the switch. For single-argument functions the one-shot result is identical
to the old loop, so the generalization is safe.
\end{gotcha}

\subsection{Formal functions and the rename dance}\label{sec:formal}

A model author often does not want to commit to a closed form for a transfer
function. Instead they declare $\phi$ as a \term{formal function}: a Sage symbolic
function left unevaluated, built with \code{function('phi\_1')(v)}. A formal
function stays inert through symbolic manipulation --- you can differentiate it,
multiply it, Taylor-expand it, all abstractly. That is exactly what we want: we
expand $\phi$ without ever saying what $\phi$ is.

The catch is what Sage produces when it Taylor-expands a formal function. The
$k$-th derivative at $0$ comes out as an atom like \code{D[0](phi\_1)(0)} (first
derivative of \code{phi\_1} at $0$) or \code{D[0,0](phi\_1)(0)} (second
derivative). These are not valid polynomial-ring coefficients. So step~6 renames
them to clean symbols:
\[
  \texttt{phi\_1(0)} \;\to\; \texttt{phi0\_1}, \qquad
  \texttt{D[0](phi\_1)(0)} \;\to\; \texttt{phi1\_1}, \qquad
  \texttt{D[0,0](phi\_1)(0)} \;\to\; \texttt{phi2\_1}.
\]
For a single-argument function the suffix is just the derivative order $k$; for a
multi-argument function it is the concatenated multi-index, so $f^{(2,1)}$ becomes
\code{f21\_<i>}. That suffix encoding is a one-liner:

\begin{lstlisting}
# field_theory.py:225
def _multi_index_suffix(multi_idx):
    return ''.join(str(k) for k in multi_idx)
\end{lstlisting}

The rename map itself, \code{ns.\_deriv\_rename\_subs}, is built during
\code{\_build\_namespace}. To know \emph{which} derivative atoms to expect, the
framework enumerates every multi-index up to the Taylor order with a small
generator:

\begin{lstlisting}
# field_theory.py:202
def _iter_multi_indices(n_args, max_total):
    if n_args == 0:
        yield ()
        return
    from itertools import product
    for combo in product(range(max_total + 1), repeat=n_args):
        if sum(combo) <= max_total:
            yield combo
\end{lstlisting}

\begin{note}
Why \code{itertools.product} rather than a tidy self-recursive generator? Because
this code runs in Jupyter under \code{\%autoreload}. A self-recursive generator
holds a reference to itself through the module's global namespace; when
\code{\%autoreload} hot-swaps the module, that self-reference is broken and the
generator fails mid-iteration. \code{itertools.product} carries no such
self-reference, so it survives the reload. The comment at
\file{field\_theory.py:210--215} records this hard-won lesson. You will see the same
defensive pattern in a few other places in the codebase.
\end{note}

Once the abstract symbols are in place, a model gives them concrete values through
its \code{specializations} lambda (step~7). For the Hawkes transfer function
$\phi=a v^2$, the specialization (at \file{hawkes\_quad\_expg.py:186}) sets
\code{phi1\_i}$\to 2a\,v^{*}_i$, \code{phi2\_i}$\to 2a$, and every higher
derivative to $0$ --- and just like that the formal $\phi$ collapses to its
quadratic reality. The \code{'expression'} authoring path in
\code{\_build\_namespace} automates even this: the model supplies the closed-form
expression, and the framework differentiates it symbolically at the all-zero point
with \code{diff} to compute each $\partial^{\alpha}f(0)$ itself, skipping any
derivative that comes out purely numeric. The legacy \code{'taylor\_coeffs'} path
instead takes a literal list of derivative coefficients and feeds them to
\code{\_poly\_taylor} (\file{field\_theory.py:193}), which assembles
$\sum_n (c_n/n!)\,x^n$ using the exact rational $1/n!$ from \code{QQ}.

\section{The convolution operator}\label{sec:conv}

This section explains step~3 of \code{expand}, and it is the one piece of the
pipeline whose \emph{whole reason to exist} is a subtle physics correctness issue.

\subsection{Why \texorpdfstring{$g \ast n$}{g*n} is not multiplication}

When a synaptic kernel $g$ multiplies a field $n$ in the action, the \code{*} in
the source text \emph{reads} as scalar multiplication but \emph{means} convolution
in time:
\[
  (g \ast n)(t) \;=\; \int g(t-s)\,n(s)\,\dd s.
\]
The framework represents this with a formal two-argument Sage function:

\begin{lstlisting}
# convolution.py:70
Conv = _sr_function('Conv', nargs=2)
\end{lstlisting}

\begin{defn}
\code{Conv(kernel, field)} is a formal time-convolution operator: by convention the
\emph{kernel} is the first argument and the \emph{field} is the second. It never
auto-evaluates --- it stays a formal symbol through all symbolic manipulation, so
each downstream consumer (the stationary mean field, the Fourier propagator, the
time-domain integrator) can apply its own reduction. The \code{nargs=2} declaration
makes Sage reject malformed calls like \code{Conv(g)} or \code{Conv(g, n, extra)}
immediately rather than producing a silently wrong expression.
\end{defn}

The asymmetry between the two arguments is the entire point. Only the second
argument (the field) carries time-domain content; the first (the kernel) is a
fixed translation-invariant filter. Consider a conductance-style term
$v\cdot\mathrm{Conv}(g,n)$ and expand both fields around the saddle,
$v=v^{*}+\delta v$, $n=n^{*}+\delta n$:
\[
  v\cdot\mathrm{Conv}(g,n)
  \;=\; v^{*}n^{*} \;+\; v^{*}\,g\,\delta n \;+\; \delta v\,n^{*} \;+\; \delta v\,g\,\delta n.
\]
Notice $\delta v$ multiplies the bare $n^{*}$ --- \emph{not} $g\cdot n^{*}$. A naive
``flatten $\mathrm{Conv}(g,n)\to g\cdot n$'' would instead give
$v\cdot g\cdot n = v^{*}g\,n^{*}+\dots$, spuriously attaching the kernel to the
$\delta v$ side of the bilinear. That single wrong coefficient would poison the free
propagator. The \code{Conv} operator exists precisely so this expansion comes out
right.

\subsection{The four reduction rules}

\code{reduce\_conv\_in\_action} (\file{convolution.py:110}) resolves every
\code{Conv} atom by handing Sage a Python callback through
\code{.substitute\_function(Conv, \_reduce\_arg)} --- Sage's mechanism for replacing
every occurrence of a formal function with the result of a callback that receives
the function's arguments. The callback applies, in order:

\begin{description}
  \item[Rule 0 --- pre-Taylor.] When a \code{taylor\_order} is supplied (it always
        is, from \code{expand}), Taylor-expand the \emph{argument} first, so a
        nonlinear inner field like $\phi(v)$ becomes a polynomial that the linearity
        rule can distribute over. A no-op for an already-polynomial argument such as
        a bare \code{n[j]}.
  \item[Rule 1 --- linearity.] $\mathrm{Conv}(g, a+b)\to\mathrm{Conv}(g,a)+\mathrm{Conv}(g,b)$,
        recursing over the summands of a sum.
  \item[Rule 2 --- pull constants.] $\mathrm{Conv}(g, c\cdot X)\to c\cdot\mathrm{Conv}(g,X)$
        for any factor $c$ with no fluctuation-field dependence (time-constant
        factors come out of the time convolution).
  \item[Rule 3 --- DC reduction.] $\mathrm{Conv}(g, c)\to\hat g(0)\cdot c = c$ for a
        time-constant leaf, under the normalized-kernel assumption $\hat g(0)=1$.
  \item[Rule 4 --- defer to the Fourier transform.] $\mathrm{Conv}(g, X)\to g\cdot X$
        for a genuinely fluctuating $X$. The kernel \emph{symbol} $g$ is left in
        place; downstream it is replaced by $\hat g(\omega)$ via
        \code{model['kernel\_ft\_image']}, which is just the convolution theorem
        $(g\ast n)\hat{\phantom{x}}(\omega)=\hat g(\omega)\,\hat n(\omega)$.
\end{description}

\begin{gotcha}
The \code{normalized=True} assumption ($\hat g(0)=1$, i.e.\ every kernel integrates
to one) is hard-wired in rule~3. Non-normalized kernels would need per-kernel DC
gains, which are not yet implemented; see \file{convolution.py:182--185}. For the
neural-style kernels this framework targets --- a synaptic filter normalized to unit
area --- the assumption holds by construction.
\end{gotcha}

\subsection{Kernel attachments}

Rule~4 has a side effect that matters two phases downstream. When it emits
$g\cdot X$ it records, into the \code{attachments\_out} dictionary, which
fluctuation variables the surviving kernel symbol $g$ attached to (set-union
accumulating across reuse). \code{expand} stashes this on the namespace as
\code{ns.\_kernel\_attachments}. Later, vertex extraction reads it to recover which
\emph{physical leg} each leftover kernel symbol belongs to, so the time-domain
integrator can place that leg at the correct convolved time. We pick this thread
back up in \S\ref{sec:typed} when we meet \code{ConvVertexType}.

\section{Bigrade classification}\label{sec:bigrade}

We have reached the conceptual core. After Taylor expansion, renaming, and
specialization, the action is one big polynomial in the field generators. Step~8
turns it into something the rest of the pipeline can navigate: a dictionary keyed
by \emph{bigrade}.

\begin{defn}
The \term{bigrade} of a monomial is the pair
$(n_{\tilde{}}, n_{\mathrm{phys}})$ where $n_{\tilde{}}$ is the total exponent over
the \emph{response} (tilde) generators and $n_{\mathrm{phys}}$ is the total exponent
over the \emph{physical} (fluctuation) generators. It is the number of response legs
and the number of physical legs the corresponding diagrammatic object carries.
\end{defn}

The physical meaning of each sector is the whole reason to classify this way:

\begin{table}[h]
\centering
\caption{The bigrade sectors and their physical roles. The first three rows are
collectively the mean-field (saddle-equation) sector and must vanish at the saddle.}
\label{tab:bigrade}
\small
\begin{tabular}{@{}lll@{}}
\toprule
bigrade $(n_{\tilde{}},n_p)$ & name & meaning \\
\midrule
$(0,0)$        & constant      & must vanish (no constant piece at the saddle) \\
$(1,0)$        & tadpole       & response-linear; must vanish at the saddle \\
$(0,1)$        & EOM residual  & physical-linear; must vanish at the background \\
$(1,1)$        & \textbf{free action}   & the bilinear kernel; its inverse is the free propagator \\
$(\ge 2,\,0)$  & \textbf{noise kernel}  & response-only, $\ge$ quadratic; the noise sources \\
total $\ge 3$  & \textbf{interaction vertex} & every higher-order monomial is a vertex \\
\bottomrule
\end{tabular}
\end{table}

The mechanics are remarkably short because the polynomial ring does the heavy
lifting. First, \code{\_sr\_to\_ring} (\file{field\_theory.py:241}) converts the
\code{SR} expression into a ring element. It expands the expression, splits it into
summands (using the name-based \code{add} test from the earlier gotcha), and for
each summand reads off the degree in each field generator, stripping it away and
accumulating whatever symbolic coefficient remains:

\begin{lstlisting}
# field_theory.py:272-281  (inside _sr_to_ring)
for term in summands:
    exponents = []
    coeff = SR(term)
    for v_sr in gen_sr:
        deg = int(coeff.degree(v_sr))
        exponents.append(deg)
        if deg > 0:
            coeff = coeff.coefficient(v_sr, deg)
    result += SR(coeff) * R.monomial(*exponents)
\end{lstlisting}

\begin{gotcha}
The summand/single-term distinction in \code{\_sr\_to\_ring} is explicit for a
reason. If the expanded action is a single product --- say $\tau\cdot\mathrm{Dt}\cdot v$
--- then calling \code{.operands()} on it would iterate its \emph{factors}
($\tau$, $\mathrm{Dt}$, $v$), not its summands, and the coefficient $\tau\cdot\mathrm{Dt}$
would be torn apart and lost. So the code wraps any non-sum in a one-element list
\code{[expanded]} before iterating (\file{field\_theory.py:268--271}). This is the
same \code{add\_vararg} hazard wearing a different hat.
\end{gotcha}

Then \code{\_collect\_bigrade} (\file{field\_theory.py:570}) walks the ring
element's \code{.dict()} and buckets each monomial by its bigrade. Because the ring
generators are ordered tilde-first, the split index is exactly \code{n\_tilde}:

\begin{lstlisting}
# field_theory.py:570
def _collect_bigrade(poly, n_tilde: int) -> dict:
    R = poly.parent()
    sectors: dict = {}
    for exp_vec, coeff in poly.dict().items():
        n_t = int(sum(exp_vec[:n_tilde]))
        n_p = int(sum(exp_vec[n_tilde:]))
        key = (n_t, n_p)
        sectors[key] = sectors.get(key, R.zero()) + SR(coeff) * R.monomial(*exp_vec)
    return sectors
\end{lstlisting}

That five-line loop is the payoff of every design decision so far. The result is
\code{ft.\_by\_tp}, a dictionary \code{\{(n\_tilde, n\_phys): polynomial\}}. The
accessors on \code{FieldTheory} read directly off it:

\begin{itemize}
  \item \code{free\_action()} (\file{field\_theory.py:1010}) returns the $(1,1)$
        sector --- the bilinear kernel whose inverse is the free propagator
        (Chapter~\ref{ch:propagator}).
  \item \code{noise\_kernel()} (\file{field\_theory.py:1015}) returns all
        $(\ge 2,0)$ sectors --- the noise sources.
  \item \code{vertices()} (\file{field\_theory.py:1021}) returns all
        total-degree-$\ge 3$ sectors --- the interaction vertices.
  \item \code{sectors()} (\file{field\_theory.py:1027}) returns the full nonzero
        bigrade dictionary; \code{summary()} pretty-prints each labeled sector.
\end{itemize}

The human-readable labels come from the static helper \code{\_sector\_label}
(\file{field\_theory.py:1380}), which is itself a compact transcription of
Table~\ref{tab:bigrade}.

\begin{example}
\textbf{The OU quartic, classified.} With $\phi=\varepsilon x^3$ and white noise
$-D\tilde x^2$, expanding the action and classifying gives a tiny \code{\_by\_tp}.
The $(1,1)$ sector is $\tilde x\,(\partial_t+\mu)x$ --- the free OU kernel. The
$(2,0)$ sector is $-D\tilde x^2$ --- the single noise source. And the
total-degree-4 $(1,3)$ sector is $\varepsilon\,\tilde x\,x^3$ --- the lone
interaction vertex, with one response leg and three physical legs. That one vertex
is the source of every loop diagram in the quickstart's $C_{xx}(\tau)$ calculation.
\end{example}

\section{Mean-field sector verification}\label{sec:mf}

Three of the bigrade sectors --- $(0,0)$, $(1,0)$, $(0,1)$ --- are special. They are
the constant, the response-linear ``tadpole,'' and the physical-linear
``equation-of-motion residual.'' Collectively they are the \term{mean-field (saddle)
sector}, and at the correct saddle point they must \emph{vanish}. The framework does
not solve the saddle equations itself --- that is a downstream module --- but it does
two things here: it \emph{verifies} this sector vanishes under the model's declared
saddle substitutions, and then it forcibly \emph{zeroes} it.

\subsection{Why classify before substituting}

This is the single most important ordering constraint in \code{expand}, and it is
easy to get backwards. The mean-field substitutions are rewrites like
$v^{*}\to (E+\sum_j w_{ij}g\cdot n^{*}_j)/(1+\tau\,n^{*})$. If you applied them to the
\emph{whole} action before classifying, that rewrite would fire \emph{everywhere}
$v^{*}$ appears --- including inside the $(1,1)$ propagator kernel and the $\ge 2$
interaction vertices, where $v^{*}$ and $n^{*}$ must remain \emph{free symbolic
parameters} to be bound numerically much later. You would inject saddle-equation
algebra into the propagator. So the code coerces and classifies first
(\file{field\_theory.py:921--922}), and applies the saddle substitutions \emph{only}
to the three mean-field sectors, purely as a test. The long comment at
\file{field\_theory.py:913--922} states this explicitly.

\subsection{The two-tier check}

\code{\_verify\_and\_zero\_mf\_sector} (\file{field\_theory.py:585}) applies the
saddle substitutions to each monomial coefficient in the three sectors and checks it
vanishes, in two tiers:

\begin{enumerate}
  \item \textbf{Symbolic.} Substitute the mean-field conditions (then the
        specializations), and if \code{.simplify\_full() == 0}, the term passes.
  \item \textbf{Numerical fallback.} Only if the symbolic check fails, evaluate the
        residual numerically at a representative parameter point via
        \code{\_mf\_numerical\_residual}; pass if its magnitude is below a tolerance
        (default $10^{-9}$), recording a \emph{soft pass} (warned, not failed).
\end{enumerate}

Any term that passes neither tier raises an \code{AssertionError} naming the
offending bigrade, its post-substitution symbolic form, and its numerical residual.
The error message is blunt: \emph{``Either the MF solver is wrong or the action's
bigrade-$\le 1$ sector is not the saddle-eq sector.''} This assertion is a genuine
structural safety net --- it catches a miswired mean-field solver or a mis-specified
action before any diagrams are computed. After the checks, each mean-field sector
key is set to \code{R.zero()}.

\begin{note}
The numerical fallback is non-trivial machinery in its own right.
\code{\_mf\_numerical\_residual} (\file{field\_theory.py:662}) builds a representative
parameter point (preferring each parameter's declared \code{default}, since theories
ship defaults that admit a well-behaved saddle), then prefers the differential-algebraic
mean-field solver when the model declares \code{equations} --- the legacy iterative
solver cannot handle self-referential implicit saddles like
$x^{*}=-\varepsilon (x^{*})^3$ --- and falls back to the simple solver otherwise. The
companion \code{\_default\_fundamental\_point} (\file{field\_theory.py:730}) warns in
its docstring that an all-ones fallback is a poor universal default for a nonlinear
closure (for $\phi=av^2$ with $E=a=w=1$ the saddle equation $2v^2-v+1=0$ has no real
root), which is exactly why declared defaults are preferred.
\end{note}

\begin{gotcha}
Specializations run \emph{twice}. Once after the Taylor (step~7, on every sector),
and again \emph{inside} the verifier (\file{field\_theory.py:949}). The reason: the
\code{mf\_bg} lambda builds its right-hand sides \emph{before} specializations have
run, so a substitution like $v^{*}\to(\dots\phi_0\dots)$ reintroduces raw
$\phi$-Taylor tokens that need a second specialization pass to resolve.
\end{gotcha}

After verification the action is rebuilt from the now-zeroed \code{by\_tp} into
\code{self.\_S\_raw}, and \code{self.\_by\_tp} is the authoritative output. The quick
public sanity check \code{sanity\_check()} (\file{field\_theory.py:966}) simply
confirms those three sectors are \code{R.zero()} and prints a PASS/FAIL table --- and
it always prints a FAIL, even when called with \code{verbose=False}, so a failure is
never silent.

\section{Correlated-noise cumulants}\label{sec:cumulant}

This section covers step~4, which is a no-op for most models and so can be skipped on
a first read. White noise (the OU model) contributes a single local term
$-D\tilde x^2$ that lives entirely in the $(2,0)$ sector; nothing special is needed.
But \emph{colored} or \emph{correlated} external noise contributes a whole series to
the action: the cumulant-generating functional $-W_m[\tilde m]$,
\[
  S_{\mathrm{cum}} \;=\; -\sum_n \frac{1}{n!}\int \dd t_1\cdots\dd t_n
     \sum_{i_1\cdots i_n} \kappa^{(n)}_{i_1\cdots i_n}(\tau\text{'s})\,
     \tilde m_{i_1}(t_1)\cdots\tilde m_{i_n}(t_n),
\]
where $\kappa^{(n)}$ is the $n$-th cumulant of the noise and the $\tilde m$ are the
response legs the noise couples to. \code{\_build\_cumulant\_action}
(\file{field\_theory.py:293}) injects this series symbolically from the declarative
\code{model['correlated\_noises']} block, returning \code{SR(0)} when there is no such
block.

The key idea is that each per-leg kernel $K(\tau)$ is split into a
\emph{local (delta-correlated)} part and a \emph{smooth non-local} part:
\[
  K(\tau) \;=\; c_{\mathrm{local}}\,\delta(\tau) \;+\; K_{\mathrm{smooth}}(\tau).
\]
The framework peels off the local part by repeatedly taking the coefficient of
$\delta(\tau_k)$:

\begin{lstlisting}
# field_theory.py:460-463  (inside _build_cumulant_action)
c_local = K
for tau_k in tau_syms:
    try:
        c_local = c_local.coefficient(dirac_delta(tau_k))
\end{lstlisting}

The local part collapses every time integral and is injected as an ordinary monomial
(using the \code{is\_trivial\_zero()} test from the earlier gotcha, at
\file{field\_theory.py:490}). The smooth residual, by contrast, cannot collapse: it
introduces a placeholder symbol \code{z\_kappa\_...} (the \code{z} prefix again for
canonical ordering) and registers the actual kernel function in
\code{ns.\_cumulant\_kernels} for the downstream integrator to pick up. The injector
also handles \emph{cross-field} cumulants by enumerating distinct permutations of the
leg-field tuple --- so a noise correlating two different fields $[x_t,y_t]$ sums over
both orderings.

\begin{gotcha}
The smooth (non-$\delta$) residual is only handled at cumulant order $n=2$. At order
$n\ge 3$, a smooth residual is \emph{silently dropped} with a \code{warnings.warn}
(the local $\delta$ part of the same kernel, if any, is still injected correctly).
The reason is purely a limitation of the time-domain integrator downstream, which
currently allocates one relative-time variable per noise vertex and would need an
$n$-leg time map to handle higher non-local cumulants (\file{field\_theory.py:550--565}).
For the Gaussian-and-shifts noise models this machinery was built for, all
order-$\ge 3$ cumulants happen to be fully local, so the dropped branch never fires
in practice --- but it is a real ceiling for future non-local higher-order kernels.
\end{gotcha}

\section{Typed vertices and sources}\label{sec:typed}

We now cross from \file{field\_theory.py} into \file{vertices.py}. The bigrade
sectors are polynomials; the diagram machinery wants \emph{individual typed records},
one per monomial, each tagged with its legs. That conversion is what
\file{vertices.py} does.

\subsection{Leg tuples and the record types}

A \term{leg} is one attachment point of a vertex or source, tagged
\code{(field\_base\_name, population\_index)} --- for example \code{('dv', 1)} for the
first population's $\delta v$ field. The population index is 1-based, or $0$ when the
field is not indexed. A monomial's exponent in a given generator becomes that many
repeated leg entries.

The records themselves are lightweight classes that use \code{\_\_slots\_\_} (a
Python optimization that removes the per-instance \code{\_\_dict\_\_}, saving memory
and forbidding stray attributes). There are four:

\begin{description}
  \item[\code{VertexType}] (\file{vertices.py:80}) --- one interaction monomial
        (total degree $\ge 3$). Carries \code{coefficient} (the SR coupling times
        combinatorial prefactor), \code{response\_legs}, \code{physical\_legs}, and
        \code{bigrade}. Its properties \code{in\_degree} (number of physical legs),
        \code{out\_degree} (number of response legs), and \code{total\_degree} are
        the diagram-facing vocabulary.
  \item[\code{ConvVertexType}] (\file{vertices.py:134}) --- a \code{VertexType}
        subclass for a conductance-style vertex, where one or more physical legs sit
        at independent times linked to the vertex by a synaptic kernel $g(\tau)$. It
        adds \code{kernel\_attachments}: a list of dicts, each recording the kernel
        \code{'symbol'}, the \code{'leg'} it attaches to, that leg's 0-based
        \code{'leg\_index'}, and a \code{'kernel\_td\_fn'} callable
        $\tau\mapsto g(\tau)$. This is where the kernel attachments from
        \S\ref{sec:conv} are finally consumed.
  \item[\code{SourceType}] (\file{vertices.py:216}) --- one noise-kernel monomial
        ($n_{\tilde{}}\ge 2$, $n_{\mathrm{phys}}=0$). Carries \code{coefficient},
        \code{response\_legs}, \code{bigrade}, and an \code{out\_degree} property.
  \item[\code{NoiseSourceType}] (\file{vertices.py:256}) --- a \code{SourceType}
        subclass backed by a non-local cumulant kernel $\kappa^{(n)}$, whose response
        legs sit at independent times. It adds \code{cumulant\_specs}, the records
        that tie each placeholder symbol back to its kernel function. Locally
        correlated ($\delta$) cumulants stay plain \code{SourceType}.
\end{description}

The bridge from polynomial to records is \code{decompose\_sector}
(\file{vertices.py:349}). For every \code{(exponent\_vector, coefficient)} in a
sector's \code{.dict()}, it parses each nonzero generator's name into a leg tuple and
appends it (to the response legs if the generator index is below \code{n\_tilde}, else
to the physical legs), repeated by the exponent. If there are no physical legs it
emits a \code{SourceType}; otherwise a \code{VertexType}:

\begin{lstlisting}
# vertices.py:386-393  (inside decompose_sector)
n_t = len(resp_legs)
n_p = len(phys_legs)
bigrade = (n_t, n_p)
if n_p == 0:
    results.append(SourceType(SR(coeff), resp_legs, bigrade))
else:
    results.append(VertexType(SR(coeff), resp_legs, phys_legs, bigrade))
\end{lstlisting}

The generator-name parsing is done by \code{\_parse\_field\_name}
(\file{vertices.py:327}): it walks backward over trailing digits to split a name like
\code{dn2} into base \code{dn} and population index $2$ (returning index $0$ when
there are no trailing digits).

The two top-level extractors \code{extract\_vertex\_types} (\file{vertices.py:500})
and \code{extract\_source\_types} (\file{vertices.py:629}) wrap
\code{decompose\_sector} and then do the kernel/cumulant upgrade: scanning each
coefficient for kernel symbols, resolving which leg each attaches to, and promoting
the plain record to a \code{ConvVertexType} or \code{NoiseSourceType} when an
attachment resolves.

\begin{gotcha}
If a kernel symbol cannot be resolved to a specific leg, \code{extract\_vertex\_types}
deliberately \emph{leaves it in the coefficient} and emits a plain \code{VertexType}
(\file{vertices.py:589--595}). The point is fail-loud: a surviving kernel symbol in a
coefficient will trip a downstream diagnostic, which is far better than silently
treating the unresolved attachment as zero.
\end{gotcha}

\subsection{What the enumerator actually asks for}

The diagram enumerator does not consume the full vertex objects --- it only needs to
know which vertex \emph{shapes} exist. That is the entire job of
\code{available\_degrees} (\file{vertices.py:777}):

\begin{lstlisting}
# vertices.py:777
def available_degrees(vertex_types, source_types):
    interaction_degrees = {(vt.in_degree, vt.out_degree) for vt in vertex_types}
    source_degrees      = {st.out_degree for st in source_types}
    return interaction_degrees, source_degrees
\end{lstlisting}

It returns a set of \code{(in\_degree, out\_degree)} pairs from the vertices and a set
of \code{out\_degree} integers from the sources. The diagram enumerator
(\file{msrjd/enumeration/degree\_scan.py}) reads exactly this to know which
``shapes'' of vertex and source it may attach when it builds graphs --- the topic of
Chapter~\ref{ch:enumeration}.

\section{Picklability and the process boundary}\label{sec:pickle}

A short but load-bearing section, because it explains design choices in
\file{vertices.py} that otherwise look arbitrary. The downstream time-domain
integrator (Phase~J) can farm diagram evaluation across a
\code{multiprocessing.Pool}, which means the vertex and source objects may be
\emph{pickled} and shipped to worker processes. This subsystem never imports
\code{multiprocessing}, but it is engineered around that possibility. Three
consequences:

\begin{itemize}
  \item \code{\_NamespaceBoundKernel} (\file{vertices.py:25}) is a \emph{module-level}
        class, not a closure, so \code{pickle} can locate it by its fully-qualified
        name. It \emph{pre-evaluates} the user's kernel function once against the
        namespace at construction time and caches the resulting \code{SR} expression
        (\file{vertices.py:56}) --- because the full namespace is unpicklable (it
        holds inner closures), but a plain \code{SR} expression pickles cleanly.
  \item The record classes implement explicit
        \code{\_\_getstate\_\_}/\code{\_\_setstate\_\_} because \code{\_\_slots\_\_}
        classes have no \code{\_\_dict\_\_} for pickle to grab automatically.
  \item In a \code{\_\_slots\_\_} inheritance chain, the base class's pickle methods
        must reference the base \code{\_\_slots\_\_} \emph{by name}.
        \code{SourceType.\_\_getstate\_\_} therefore writes
        \code{SourceType.\_\_slots\_\_} explicitly, not \code{self.\_\_slots\_\_} ---
        on a \code{NoiseSourceType} instance the latter resolves only to the
        subclass's own \code{('cumulant\_specs',)}, which would drop the base fields
        from the pickle (\file{vertices.py:235--238}).
\end{itemize}

\begin{gotcha}
Because \code{\_NamespaceBoundKernel} pre-evaluates \emph{eagerly} at extraction time,
a kernel function that errors will throw during \code{extract\_source\_types}, not
later during integration. That is usually a feature (you find out early), but it can
surprise you if you expected lazy evaluation.

A project-wide caveat, recorded in the framework's memory: fork-based multiprocessing
has crashed this user's macOS machine. The spatial path is serial-only and the
temporal path is fork-guarded. The picklability machinery here is what \emph{allows}
these objects to cross a (guarded) process boundary at all, but the guard itself lives
elsewhere in the pipeline.
\end{gotcha}

\section{Serialization}

The last file, \file{serialize.py}, persists an expanded theory so the expensive
multivariate Taylor step is not repeated every run. \code{save\_theory}
(\file{serialize.py:52}) writes a directory with two artifacts:

\begin{itemize}
  \item \file{metadata.json} --- plain-Python metadata: format and Sage version, a
        timestamp, the model identity (\code{model\_name}, \code{model\_file},
        \code{model\_var\_name}), the expansion info (\code{taylor\_order},
        \code{n\_tilde}, \code{ring\_var\_names}), the field/parameter/kernel specs
        with callables stripped, and a sorted list of nonzero bigrade keys for quick
        inspection.
  \item \file{symbolic\_data.sobj} --- a Sage \emph{binary} object dump (via
        \code{sage\_save}) of the polynomial ring, the raw action, the bigrade
        dictionary, \code{n\_tilde}, and the propagator matrices. Sage's \code{.sobj}
        format can pickle arbitrary Sage objects --- polynomial rings, symbolic
        expressions, matrices --- that plain \code{pickle}/\code{json} cannot.
\end{itemize}

\begin{gotcha}
Two deliberate omissions. First, the model dict itself (with its lambdas) is
\emph{not} serialized; instead the metadata records the model's file path and
variable name so it can be re-imported. Second, the \code{VertexType}/\code{SourceType}
objects are \emph{not} saved either --- they are cheap to re-extract from the bigrade
dictionary on load, and they carry namespace-derived callables that are better
reconstructed than persisted.
\end{gotcha}

\code{load\_theory} (\file{serialize.py:156}) returns the \code{(metadata, data)} pair,
and \code{reload\_model} (\file{serialize.py:192}) re-imports the model dict from the
stored path. The diagram enumerator uses exactly this round-trip to re-expand a saved
theory at a \emph{higher} Taylor order on demand, when a particular diagram needs more
vertex legs than the cached expansion provides.

\begin{gotcha}
A trap for the Sage-literate. The comment at \file{serialize.py:232} says ``Use
SageMath's \code{load()} which handles .py files,'' but the body actually uses a plain
Python \code{exec(compile(code, full\_path, 'exec'), ns)}. The comment is stale. The
practical consequence: this \code{exec} does \emph{not} run Sage's preparser, so a
model file relying on Sage-only syntax (\code{\^{}} for power, a bare \code{2/3} meaning
an exact rational) would not be preparsed. Every model in this repository imports from
\code{sage.all} and uses plain Python syntax, so it works --- but do not author a model
file in Sage-preparser dialect and expect \code{reload\_model} to cope.
\end{gotcha}

\section{Where this sits in the pipeline}

Step back and place the whole subsystem. It is the front end --- phase \code{[1/7]}
--- and its job is one transform:
\[
  \underbrace{\text{model dict}}_{\text{declarative}}
  \;\xrightarrow{\ \text{\code{FieldTheory(...).expand()}}\ }\;
  \underbrace{\code{ft.\_by\_tp}}_{\{(n_{\tilde{}},n_p)\,:\,\text{polynomial}\}}.
\]

What \emph{feeds} it is the model dict, authored via the theory builders of
Chapter~\ref{ch:theory-spec}. The main driver \file{pipeline/compute.py} makes the
entry call \code{FieldTheory(model, taylor\_order).expand()} as its first phase.

What \emph{consumes} it splits three ways, all reading off \code{ft.\_by\_tp}:

\begin{itemize}
  \item \code{ft.free\_action()} (the $(1,1)$ sector) feeds the \textbf{propagator
        extractor} of Chapter~\ref{ch:propagator}, which inverts the bilinear kernel
        to get the free propagator $\Gret$.
  \item \code{extract\_source\_types(ft)} and \code{extract\_vertex\_types(ft)} feed
        \code{available\_degrees}, which feeds the \textbf{diagram enumerator} of
        Chapter~\ref{ch:enumeration}.
  \item The full \code{VertexType}/\code{SourceType}/\code{ConvVertexType}/\code{NoiseSourceType}
        records feed the \textbf{integrators} (temporal and spatial) of Parts~IV
        and~V, which finally turn the diagrams into numbers.
\end{itemize}

Everything else in this manual builds on the classified action you have just watched
get constructed. With the free action in hand, the next chapter inverts it to obtain
the propagator.
